\documentclass[11pt]{article}
\usepackage[margin=1in]{geometry}
\usepackage[T1]{fontenc}
\usepackage[utf8]{inputenc}
\usepackage{lmodern}
\usepackage{microtype}
\usepackage{amsmath,amssymb}
\usepackage{booktabs}
\usepackage{graphicx}
\usepackage{hyperref}
\usepackage{enumitem}
\usepackage{tabularx}

\title{Math4AI Capstone Report Title}
\author{Team Member 1 \and Team Member 2 \and Team Member 3}
\date{}

\begin{document}
\maketitle

\begin{abstract}
State the central question, the models compared, the main experimental evidence, and the conclusion in 120--180 words.
\end{abstract}

\section{Introduction}
Frame the capstone question and explain why comparing a linear model to a one-hidden-layer neural network is scientifically meaningful.

\section{Background and Mathematical Setup}
Define the AI concepts used in the project and summarize the required mathematical setup. Explain why softmax regression is a probabilistic model, why cross-entropy is negative log-likelihood, and what the hidden layer changes geometrically.

\section{Methods}
Describe the datasets, preprocessing, softmax regression, neural network, training protocol, and evaluation metrics. Make the digits protocol explicit and state exactly how validation and test data were used.

\section{Experiments}
Report the required core comparisons and ablations. Prefer a small number of well-designed figures over many shallow plots. Explain where the linear model is sufficient, where the hidden layer changes the geometry, and what evidence supports each conclusion.

\section{Advanced Track}
Include exactly one of the following:
\begin{itemize}[leftmargin=1.4em]
    \item Track A: PCA/SVD and input geometry
    \item Track B: Prediction confidence and reliability
\end{itemize}

\section{Discussion}
Interpret the results. State clearly where linear structure was sufficient, where nonlinearity helped, and why. Use consistent notation and define symbols before using them.

\section{Limitations}
Explain what the project does not test and which conclusions should not be overgeneralized.

\section*{References}
Use a consistent citation style.

\end{document}
